% Generated draft: requires mathematical and source review.
\documentclass[11pt]{article}
\usepackage[margin=25mm]{geometry}
\usepackage{fontspec}
\setmainfont{Latin Modern Roman}
\usepackage{amsmath,amssymb,mathtools,unicode-math}
\setmathfont{Latin Modern Math}
\usepackage{enumitem,needspace,xurl,hyperref}
\setlength{\parindent}{0pt}
\setlength{\parskip}{5pt}
\setlength{\emergencystretch}{3em}
\newcommand{\N}{\mathbb{N}}
\newcommand{\Z}{\mathbb{Z}}
\newcommand{\Q}{\mathbb{Q}}
\newcommand{\R}{\mathbb{R}}
\newcommand{\C}{\mathbb{C}}
\newcommand{\F}{\mathbb{F}}
\newcommand{\A}{\mathbb{A}}
\newcommand{\PP}{\mathbb{P}}
\newcommand{\E}{\mathbb{E}}
\newcommand{\eps}{\varepsilon}
\newcommand{\dd}{\,\mathrm{d}}
\begin{document}
\section*{UnsolvedMath ID: 700030}
\textbf{Source problem code:} catalogue entry
\subsection{Definitions and mathematical statement}
Let \(G\) be a group. A partial order on \(G\) is a reflexive, antisymmetric, transitive relation on \(G\); it is a linear order if any two elements are comparable. For \(C\in\{L,R,B\}\), call a partial order \(C\)-invariant if it satisfies the corresponding condition:
\[
\begin{aligned}
L:&\quad x\leq y\Longrightarrow gx\leq gy,\\
R:&\quad x\leq y\Longrightarrow xg\leq yg,\\
B:&\quad x\leq y\Longrightarrow gxh\leq gyh,
\end{aligned}
\qquad(x,y,g,h\in G).
\]
For a fixed convention \(C\), say that \(G\) is \emph{pro-orderable under convention \(C\)} if every \(C\)-invariant partial order \(\leq\) on \(G\) admits a \(C\)-invariant linear order \(\preccurlyeq\) on \(G\) extending it, meaning
\[
\forall x,y\in G,\qquad x\leq y\Longrightarrow x\preccurlyeq y.
\]
A group \(G\) is \emph{simple} if \(G\neq\{1\}\) and its only normal subgroups are \(\{1\}\) and \(G\).

\textbf{Part (c), attributed to A. I. Mal'cev and L. Fuchs.} Does there exist a simple pro-orderable group? For each fixed convention \(C\), the corresponding precise existence question is
\[
\exists G\;\bigl(G\text{ is simple and pro-orderable under convention }C\bigr).
\]
The supplied source does not specify whether its group orderings are left invariant, right invariant, or invariant on both sides. Thus the displayed family makes this unresolved qualification explicit; it does not identify which member is the intended source question. Reading ``ordering'' as an arbitrary order on the underlying set would omit a potentially essential group compatibility requirement.

\textbf{Archived parts retained from the source.} Under the same unresolved ordering convention, part (a) asks whether the wreath product of every pair of pro-orderable groups is pro-orderable, and part (b), attributed to A. I. Mal'cev, asks whether every subgroup of every pro-orderable group is pro-orderable. The supplied excerpt does not specify whether the wreath product in (a) is restricted or unrestricted. Its archive reports a negative answer to each of (a) and (b), attributed to V. M. Kopytov. This is a supplied historical report, not an independently verified assertion here; it does not supply an answer to (c).
\subsection{Short English statement}
Part (c) asks whether a nontrivial simple group can have the property that every compatible partial order extends to a compatible linear order. The required compatibility convention is absent from the supplied excerpt and needs clarification. The record labels the overall problem partially solved and reports negative answers to archived parts (a) and (b), while retaining (c) in the main section.
\subsection{Sources}
\begin{itemize}
\item Supplied local UnsolvedMath registry record, permanent ID \(700030\), exact source problem code \texttt{catalogue entry}, title \emph{Kourovka Notebook Problem 1.35}; upstream revision \texttt{d371ee917ac9647fbb06cd6fc529deec6f3dda82}.
\item \emph{The Kourovka Notebook (Unsolved Problems in Group Theory)}, Problem 1.35, first issue (1965), described in the supplied record as updated in the twenty-first issue (2026), \texttt{arXiv:1401.0300v48}. The collection description names E. I. Khukhro and V. D. Mazurov as editors. The overall problem is attributed to M. I. Kargapolov; part (c) to A. I. Mal'cev and L. Fuchs; part (b) to A. I. Mal'cev.
\item Supplied Archive of Solved Problems attribution for the negative answers to parts (a) and (b): V. M. Kopytov, \emph{Algebra i Logika}, volume 5, no. 6 (1966), pp. 27--31, in Russian. No article title or direct article link is supplied.
\item The supplied status review is dated \texttt{2026-10-16} and labels the record \texttt{partially\_solved}, reporting archived solved parts and remaining parts in the main section. This is supplied dated curation, not an independently verified literature assessment. No external source was visited for this draft.
\end{itemize}
\end{document}
